\chapter{Introduction and physical motivation}
\label{ch:intro}

\section{Problem statement}
Consider steady laminar forced convection in a two-dimensional slit channel. A finite segment of the
upper wall is not rigid: it consists of a thin thermoelastic layer whose outer face is heated by a
uniformly applied heat flux, as realised, for example, by a film heater. Conduction across the layer
produces a through-thickness temperature difference whose magnitude is governed by the wall
conductivity and thickness. Linear thermoelasticity associates with this gradient a free thermal
curvature; with the segment ends pinned, the segment bows toward the interior of the passage. The
resulting displacement redistributes the flow cross-section, alters the near-wall velocity and
temperature gradients, and therefore changes the interfacial convective heat flux. That flux, in
turn, is one of the boundary data of the conduction problem in the wall. The wall temperature, the
deformation, and the fluid thermal field are thus coupled in both directions, and a consistent
solution requires their simultaneous determination.

The present document formulates this coupled problem mathematically, derives the reduction implied
by plane strain and by the thin-wall geometry, evaluates a baseline operating point analytically,
and specifies a numerical formulation and a verification hierarchy for its eventual computation. No
simulation results are reported; the distinction between established theory, results derived here
for this configuration, and quantities that remain to be computed is maintained throughout.

\section{Physical motivation}
Wall-guided micro heat exchangers and microreactors are increasingly fabricated from polymers and
polymer composites, which combine low thermal conductivity with a comparatively high coefficient of
thermal expansion. For the same imposed heat flux, such walls sustain through-thickness temperature
drops orders of magnitude larger than metallic walls of equal thickness, and their free thermal
curvature is correspondingly larger (the comparison is quantified in
Chapter~\ref{ch:baseline}). When the wall thickness is comparable to the channel height, the
associated deformation is not a perturbation of the passage geometry but a modification of it, and
heat-transfer predictions that assume a fixed geometry become geometry-dependent in a
thermally self-consistent way.

Deformation of flow passages by thermal loading is well known from thermoelastic-contact and
regenerative-cooling analyses, in which the deformation is typically computed a posteriori from a
frozen flow solution, i.e.\ in one direction of coupling. In the configuration considered here, the
deformation is itself the principal quantity of interest and feeds back on the convective boundary
condition that produces it. A number of natural modelling questions follow and motivate the present
work: under which combinations of wall material, geometry, and hydraulic operation is the thermal
bowing comparable to the channel height; what bounds does thermal buckling of the pinned segment
impose on admissible heating; how does the interaction behave under imposed flow rate relative to
imposed pressure drop; and how large are the residual stresses carried by the wall.

\section{Scope of the present study}
The analysis is confined to a two-dimensional geometry, steady state, laminar incompressible flow,
and linear isotropic thermoelasticity under plane strain, with the wall equation extended to the
pinned beam--column form at finite membrane force. The baseline treats imposed heat flux on
the outer face of the deformable segment; a variant with imposed outer-face temperature is retained
as a separately documented operating regime. The solid is modelled by two-dimensional elasticity and
steady conduction; the reduced wall equations are used for design estimates
and benchmark comparisons. Buoyancy is disposed of at the baseline by a length-consistent
Richardson-number estimate (three orders of magnitude below the mixed-convection range), with a
regime guard that re-opens the question should a parameter study drive the number upward, as
discussed in Chapters~\ref{ch:baseline} and~\ref{ch:limitations}.

Transient coupling, mesh dynamics, and the associated discrete geometric conservation-law
requirements \cite{lesoinne1996geometric} are outside the present scope and are recorded as the
conditions of a possible start-up study. Likewise, creep of the polymer wall, manufacturing
tolerances, and contact resistance at the interface are excluded from the baseline and are listed
among the model limitations.

\section{Relation to existing work}
Fluid--structure interaction with heat transfer has been studied for cavity and channel geometries
in which the deformation is driven by fluid traction \cite{soti2015steady,alsabery2019numerical,%
saleh2025laminar}, and fully nonlinear thermo-fluid--structure interaction of flexible cantilevers
in a differentially heated cavity has been computed \cite{kumar2026thermo}.
In the configuration studied here, the deformation mechanism is instead the wall's own thermal
eigen-curvature, the flow is a confined forced-convection slit at moderate Reynolds number, and the
coupling closes through the conjugate thermal boundary condition. A detailed comparison, the
resulting positioning of the intended contribution, and the limits of that positioning imposed by
the available literature search are developed in Chapter~\ref{ch:contribution}.

\section{Organisation of the document}
Chapter~\ref{ch:model} fixes the geometry, assumptions, sign conventions, and material data.
Chapter~\ref{ch:equations} states the governing equations and the constitutive relations together
with the principles from which they follow. Chapter~\ref{ch:derivations} derives the plane-strain
reduction and the analytical relations used for design estimates. Chapter~\ref{ch:interface}
assembles the boundary and interface conditions. Chapter~\ref{ch:dimensionless} identifies the
dimensionless groups and the admissibility bounds of the regime. Chapter~\ref{ch:baseline}
evaluates the baseline operating point in full arithmetic detail. Chapter~\ref{ch:numerics}
proposes the finite-element formulation and solution procedure. Chapter~\ref{ch:validation}
specifies the verification and validation methodology. Chapter~\ref{ch:contribution} reviews the
related literature and states the proposed contribution and its bounds.
Chapter~\ref{ch:limitations} collects the limitations. Appendices contain the notation, the
arithmetic calculations in independently checkable cells, a dimensionless dictionary, and the audit
record (discrepancy register, consolidated evidence status, unresolved source checks, and the
conditions attached to the staged implementation decisions).
